%%  Test/exam template at the University of Toronto.
%%  Original template (c) 2015-2022 Francois Pitt <fpitt@cs.utoronto.ca>, CC BY-SA 4.0.
%%  Adapted for MIE 1626 Test 1 (Midterm).

\RequirePackage[l2tabu,orthodox]{nag}  % warn about common LaTeX pitfalls

\documentclass[11pt,twoside]{article}

\usepackage{uttest}
\usepackage{macros-exams}
\usepackage{esc180} % shared course helpers (headers, blank page, tables)


%%%%%%%%%%%%%%%%%%%%%%%%%%%%%%%%%%%%%%%%%%%%%%%%%%%%%%%%%%%%%%
%%%%%%%%%%%%%%%%%%%%%% MODIFY THIS PART %%%%%%%%%%%%%%%%%%%%%%
%%%%%%%%%%%%%%%%%%%%%%%%%%%%%%%%%%%%%%%%%%%%%%%%%%%%%%%%%%%%%%
%\setboolean{restricted}{true}  % uncomment if exam is restricted
\newcommand*{\faculty}{Faculty of Applied Science and Engineering}
\renewcommand*\testheadverso{TEST}
\renewcommand*\testheadrecto{1}
\renewcommand*\examperiod{MAY 2026}
\newcommand*{\courseid}{MIE\,1626\,H}
\newcommand*{\coursetitle}{Data Science Methods and Statistical Learning}
\newcommand*{\calculator}{None}
\newcommand*{\examtype}{Open: reference sheet}
\newcommand*{\examiner}{M.~Guerzhoy}
\renewcommand*\duration{1 hour 10 minutes}
%%%%%%%%%%%%%%%%%%%%%%%%%%%%%%%%%%%%%%%%%%%%%%%%%%%%%%%%%%%%%%
%%%%%%%%%%%%%%%%%%%%%%%%%%%%%%%%%%%%%%%%%%%%%%%%%%%%%%%%%%%%%%
%%%%%%%%%%%%%%%%%%%%%%%%%%%%%%%%%%%%%%%%%%%%%%%%%%%%%%%%%%%%%%

\begin{document}
\thispagestyle{plain}

\ifsolutions
\ifanswerspace
\heading{Summary of Marking Guidelines:}
\begin{itemize}
\item
    If you are seeing this and not a cover page, you have `solutions' set to
    true in the exam.tex file.
\end{itemize}

\vfillstretch

\newpage
\fi%answerspace
\else%solutions
\ifanswerspace
\begin{center}
\bfseries\large
    \uppercase\expandafter{University of Toronto}\\
    \uppercase\expandafter{\faculty}\\[1.75ex]
    \uppercase\expandafter{\testhead, \examperiod}\\[1.75ex]
    \uppercase\expandafter{Duration: \duration}\\[1.75ex]
    \uppercase\expandafter{\courseid} --- {\coursetitle}\\[1.75ex]
    Calculator Type: {\calculator}\\
    Exam Type: {\examtype}\\[1.75ex]
    Aids allowed: the reference sheet distributed with the test\\
    Examiner(s): {\examiner}
\end{center}

\vfillstretch[.5]

\studentinfo

\vfillstretch[.75]

\begin{center}
\rule[.1875\baselineskip]{\textwidth}{1pt}\\
\larger[1.5]%
    \textsl{Do \strong{not} turn this page
    until you have received the signal to start.} \\
    In the meantime,
    please read the instructions below \textsl{carefully}.
\rule[.1875\baselineskip]{\textwidth}{1pt}\\
\end{center}

\vfillstretch[.5]

\noindent\parbox{5.25in}{%
    \qquad
    This test paper consists of
    \ref{LastQN} {questions} on
    \pageref{LastPage} {pages}
    (including this one), printed on both sides of the paper.
    \altemph{When you receive the signal to start,
    please make sure that your copy is complete,
    and fill in the identification section above.}

    \qquad
    Answer each question directly on this paper, in the space provided. Use the pages at the end of the test for extra space. If you use extra pages, indicate that you have done so in the space under the question.

    \qquad
    Write up your solutions carefully! Comments are \emph{not} required to
    receive full marks. However, they may help us mark your answers, and
    part marks \emph{might} be given for partial solutions with comments
    clearly indicating what the missing parts should accomplish.

    \qquad
    Use the \textbf{R} programming language. You may assume the
    \texttt{tidyverse} is loaded and that the data frames named in each
    question are already available. Verb signatures
    (\texttt{group\_by}, \texttt{summarize}, \texttt{filter},
    \texttt{lm}, \texttt{glm}, \texttt{sample}, \texttt{replicate},
    \texttt{plogis}, \dots) are on the reference sheet, so exact syntax need
    not be memorized and minor syntax slips are not penalized. You are
    \emph{not} required to write \texttt{ggplot} code.

    \qquad
    When you are asked to write code,
    \emph{no} error checking is required:
    you may assume that all input and argument values are valid,
    except where explicitly indicated otherwise.

    \qquad
    You \textbf{must} write your student number on every odd-numbered page of the test (except empty pages). \textbf{Failure to do so will result in a 3-point penalty.}
}%parbox
\hfillstretch
\markstable[\textsc{\larger[.5]Marking Guide}]
\hfillstretch[.25]

\vfillstretch
\newpage
\fi%answerspace
\fi%solution



% Question 1: Functions in R
\question{}

\subquestion[6]{}
Suppose you are given three one-argument functions \verb|f|, \verb|g|, and \verb|h| (already defined). Write R code that defines a function \verb|chained(x)| that returns \verb|f(g(h(x)))|, and show one example call.

For concreteness you may \emph{make up} simple definitions of \verb|f|, \verb|g|, and \verb|h| (for example doubling, adding one, squaring) so that your example produces a concrete number.

\ifanswerspace
\vfillstretch
\fi

\subquestion[6]{}
Consider the function

\begin{verbatim}
SpecialSquare <- function(num){
  if(num == 42){
    return(0)
  }else{
    return(num**2)
  }
}
\end{verbatim}

Write R code that computes \verb|SpecialSquare| of \strong{each} of the integers \verb|seq(1, 50)|, returning a vector of length 50. (There is more than one acceptable way to do this in R.)

\ifanswerspace
\vfillstretch
\fi

\ifanswerspace\newpage\fi

% Question 2: Gapminder
\question[8]{}

The \verb|gapminder| data frame contains one row per country per year. Its first rows look like this:

\begin{verbatim}
      country continent year lifeExp      pop gdpPercap
1 Afghanistan      Asia 1952  28.801  8425333  779.4453
2 Afghanistan      Asia 1957  30.332  9240934  820.8530
3 Afghanistan      Asia 1962  31.997 10267083  853.1007
4 Afghanistan      Asia 1967  34.020 11537966  836.1971
5 Afghanistan      Asia 1972  36.088 13079460  739.9811
6 Afghanistan      Asia 1977  38.438 14880372  786.1134
\end{verbatim}

The years recorded are \verb|1952, 1957, 1962, 1967, 1972, ...| (five-year steps), and \verb|continent| takes the values \verb|Africa, Americas, Asia, Europe, Oceania|.

Write R code that computes the \strong{continent with the highest median \texttt{lifeExp} in the year 1967}. Your code should return the name of that continent (or a one-row summary that identifies it).

\ifanswerspace
\vfillstretch
\fi

\ifanswerspace\newpage\fi

% Question 3: Baby names
\question[8]{}

The \verb|babynames| data frame has one row per (\verb|year|, \verb|sex|, \verb|name|) combination, giving the number \verb|n| of US babies of that sex given that name in that year. Its first rows:

\begin{verbatim}
  year sex      name    n       prop
1 1880   F      Mary 7065 0.07238359
2 1880   F      Anna 2604 0.02667896
3 1880   F      Emma 2003 0.02052149
4 1880   F Elizabeth 1939 0.01986579
5 1880   F    Minnie 1746 0.01788843
6 1880   F  Margaret 1578 0.01616720
\end{verbatim}

Write R code that computes \strong{the number of babies born corresponding to the most popular baby name in 1982} --- that is, find the name with the largest total count \verb|n| (summed over both sexes) among babies born in 1982, and report that total count.

\ifanswerspace
\vfillstretch
\fi

\ifanswerspace\newpage\fi

% Question 4: Estimating a slope three ways
\question{}

Data were generated from a relationship of the form
\[
  y = a\,x + \varepsilon, \qquad \varepsilon \sim \mathcal{N}(0, \sigma^2),
\]
and stored in a data frame \verb|d| with columns \verb|x| and \verb|y| (60 rows). The first rows are:

\begin{verbatim}
      x      y
1 0.675  1.234
2 0.254 -9.641
3 6.305 21.037
4 9.141  9.622
5 6.671 11.372
6 8.744 22.385
\end{verbatim}

Running \verb|summary(lm(y ~ x, data = d))| produces:

\begin{verbatim}
Call:
lm(formula = y ~ x, data = d)

Residuals:
   Min     1Q Median     3Q    Max
-9.300 -2.912 -0.212  2.360 10.367

Coefficients:
            Estimate Std. Error t value Pr(>|t|)
(Intercept)   -2.080      1.151  -1.806   0.0761 .
x              2.202      0.211  10.435 6.24e-15 ***
---
Signif. codes:  0 '***' 0.001 '**' 0.01 '*' 0.05 '.' 0.1 ' ' 1

Residual standard error: 4.472 on 58 degrees of freedom
Multiple R-squared:  0.6525,    Adjusted R-squared:  0.6465
F-statistic: 108.9 on 1 and 58 DF,  p-value: 6.245e-15
\end{verbatim}

\ifanswerspace\newpage\continuequestion\fi

\subquestion[6]{}
Using \strong{only the summary above}, give a point estimate of $a$ together with an approximate 95\% confidence interval. Show the arithmetic.

\ifanswerspace
\begin{verbatim}




\end{verbatim}
\vfillstretch
\newpage
\fi

\subquestion[6]{}
Describe how you would estimate $a$ by \strong{directly minimizing the appropriate linear-regression loss function}, without calling \verb|lm|. Write down the loss as a function of the candidate slope (and intercept), and \strong{write R code} that finds the minimizing value. \emph{(We have not covered how to quantify uncertainty in this case, so a point estimate is all that is required here.)}

\ifanswerspace
\begin{verbatim}










\end{verbatim}
\vfillstretch
\newpage
\fi

\subquestion[6]{}
Describe how you would estimate $a$ \strong{and its uncertainty using a Bayesian approach}, so that your answer corresponds to a region that contains the true value of $a$ with 95\% probability. State your prior, what the posterior is over, and how you would extract the 95\% region from posterior samples. \emph{(This was not covered directly --- some creative thinking is expected.)}

\ifanswerspace
\begin{verbatim}










\end{verbatim}
\vfillstretch
\fi

\ifanswerspace\newpage\fi

% Question 5: Weak relationship, tiny p-value
\question[12]{}

Write R code that \strong{generates data} \verb|x| and \verb|y| for which the relationship between them is \emph{weak} (low $R^2$), yet the p-value reported by \verb|lm(y ~ x)| for the slope comes out \emph{very small}. (You do not need to guarantee any exact threshold --- it just needs to be very small.)

Explain in one or two sentences \strong{why} your construction produces this combination, and what it tells us about interpreting small p-values.

\ifanswerspace
\vfillstretch
\fi

\ifanswerspace\newpage\fi

% Question 6: Insurance on the Titanic
\question{}

You have the \verb|titanic| (training) data frame --- one row per passenger --- and you are tasked with selling drowning insurance to passengers as they board. The first rows of the relevant columns are:

\begin{verbatim}
  Survived Pclass    Sex Age SibSp Parch    Fare
1        0      3   male  22     1     0  7.2500
2        1      1 female  38     1     0 71.2833
3        1      3 female  26     0     0  7.9250
4        1      1 female  35     1     0 53.1000
5        0      3   male  35     0     0  8.0500
6        0      3   male  NA     0     0  8.4583
\end{verbatim}

Here \verb|Survived| is 1 if the passenger survived and 0 otherwise; the data frame has 891 rows. You set the insurance \strong{premium} per passenger. The \strong{payout} in case of drowning is always \$100. If the passenger survives, they do \strong{not} get the premium back, and they do \strong{not} get the payout.

\subquestion[8]{}
Propose a reasonable formula representing a single passenger's \strong{decision} of whether to buy the insurance, given the premium you offer them. State your assumptions about how a passenger values the \$100 payout and the premium, and write the condition under which they buy.

\ifanswerspace
\vfillstretch
\fi

\subquestion[8]{}
Suppose you are allowed to use only \verb|Pclass| and \verb|Age| to price each passenger. State how you would produce a formula for the \strong{smallest premium that does not lead to an expected loss} for you, and state precisely how you would use that formula to set a price for a new passenger of given \verb|Pclass| and \verb|Age|. \emph{(Hint: relate the break-even premium to an estimated probability of drowning.)}

\ifanswerspace
\vfillstretch
\fi

\ifanswerspace\newpage\fi

% Question 7: Critiquing the model
\question[10]{}

In \strong{English} (no code), critique the insurance pricing model you produced in Question~6: give \strong{two} distinct reasons why the prices might be wrong if the model were actually used. Be specific about the mechanism behind each reason.

\ifanswerspace
\vfillstretch
\fi

\ifanswerspace\newpage\fi

% Question 8: The desk-drawer (file-drawer) effect
\question[16]{}

A research community has \strong{100 researchers}. Each runs one experiment. \strong{10\%} of the hypotheses being tested are \emph{true} (there is a real effect) and \strong{90\%} are \emph{false} (no real effect). A result is \emph{reported} (published) only when the researcher \strong{rejects the null hypothesis}; experiments that fail to reject the null are filed away in a desk drawer and never reported.

Using a \strong{simulation} (you must use a simulation --- not a closed-form formula), estimate \strong{what proportion of the reported results are false positives}. The output of your code should be that proportion.

Write the R code (\verb|sample|, \verb|runif|, \verb|rbinom|, \verb|replicate|, \dots{} are available) and \strong{state any assumptions} you make --- in particular the significance level $\alpha$ used for false hypotheses and the statistical power used for true hypotheses.

\ifanswerspace
\vfillstretch
\fi

\ifanswerspace\newpage\fi


\newpage
\vfillstretch
\newpage
\vfillstretch


\ifanswerspace\ifsolutions\else\newpage\blankpage\blankpage\fi\fi

%% Last Page
\ifanswerspace\ifsolutions\else
\cfoot{Total Marks = \ref{TotalMarks}}
\fi\fi
\rfoot{\scshape end of test 1}

\end{document}
